\documentclass[11pt]{article}
\usepackage[margin=1.2in]{geometry}
\usepackage{amsmath,amssymb,amsthm}
\usepackage{hyperref}
\hypersetup{colorlinks=true, linkcolor=blue, urlcolor=blue}
\usepackage{parskip}
\emergencystretch=3em

\newcommand{\R}{\mathbb{R}}
\newcommand{\code}[1]{\texttt{#1}}

\title{Tide retrospective: the correction coefficient functions\\
\large continuity and polynomial growth for the numerator stage}
\author{Claude (Anthropic), with GPT-5.6 Sol consults}
\date{2026-08-10}

\begin{document}
\maketitle

\begin{abstract}
The integrated numerator expansion's coefficients are
$a_j = \int z^\alpha e^{-T_2(z)} P_j(z)\,dz$, where $P_j(z)$ is the
graded correction coefficient of the previous stage instantiated at
the exponent data $a_s = V_s(z)$ of the exponent split. For those
integrals to exist, and for the dominated-convergence step that
integrates the pointwise expansion, the coefficient functions
$z \mapsto P_j(z)$ need exactly two properties: continuity and
polynomial growth. Both are inductions on the power through
\code{Polynomial.coeff\_mul} --- the coefficient of a product is the
antidiagonal convolution --- seeded by the if-then-else description
of the exponent polynomial's coefficients and the operator-norm
bound on diagonal Taylor terms. The growth exponent bookkeeping is
the mathematical content: a coefficient of the $i$-th power at
degree $k$ grows at most like $(1+\|z\|)^{k+2i}$, because each
exponent-polynomial coefficient at degree $u$ carries $\|z\|^{u+2}$
and the antidiagonal constraint sums the degrees to $k$.
\end{abstract}

\section{Setting}

Stage 5a of the forward-expansion programme, the consult-independent
half of the numerator stage: whatever the final integration
architecture, these functions and bounds are needed. A parallel
architecture consult on the rest of stage 5 ran during this tide and
is archived in the tide log; it confirmed the decomposition, with
the requirement (already satisfied here) that the coefficient
theorems cover arbitrary degree $j$, not only $j \le N$, because the
polynomial tail estimate consumes coefficients up to degree $N^2$.
It also confirmed the unrestricted-$x$ exponential remainder route
over a second, smaller localization window: the smaller window would
need its own tail theorems and annulus estimates, and at $N = 0$ the
Peano remainder admits no absolute bound on it anyway.

\section{The thing formalised}

Coefficient descriptions: \code{exponentPoly\_coeff} (the exponent
polynomial's coefficient at $u$ is $a_u$ on $1 \le u \le N$ and zero
otherwise, by collapsing the \code{coeff\_X\_pow} indicator sum) and
\code{gradedExpPoly\_coeff} (a finite sum over coefficients of
powers). The functions: \code{correctionCoeffFn} with
\code{correctionCoeffFn\_zero} $= 1$. Continuity:
for the exponent polynomial's coefficients the if-condition is
$z$-independent, so \code{split\_ifs} reduces to the seabed's
\code{taylorHomogeneousTerm\_continuous}; for powers, induction with
the antidiagonal convolution; for the correction coefficients, a
finite sum of the above. Growth: \code{abs\_exponentPoly\_coeff\_le}
($C_u (1+\|z\|)^{u+2}$ from the operator-norm bound),
\code{abs\_exponentPoly\_pow\_coeff\_le}
($C_{i,k}(1+\|z\|)^{k+2i}$, the induction using
\code{choose} to extract per-degree constants and the antidiagonal
membership $x_1 + x_2 = k$ to match exponents), and
\code{abs\_correctionCoeffFn\_le}
($C_j (1+\|z\|)^{j+2N}$, absorbing each $i \le N$ into the maximal
exponent since the base exceeds one).

\section{Where progress stalled}

Only small rounds. \code{positivity} cannot see the nonnegativity of
constants extracted by \code{choose} (they are opaque local
functions), so product nonnegativity must be assembled by hand with
\code{mul\_nonneg}. A $\leftarrow$\code{pow\_add} rewrite found no
adjacent power factors until \code{mul\_mul\_mul\_comm} reassociated
the product. A \code{gcongr} emitted its side goals in an unexpected
order; explicit monotonicity lemmas
(\code{pow\_le\_pow\_right}$_0$ twice composed with
\code{mul\_le\_mul\_of\_nonneg\_left}) were more predictable.

\section{Mathlib lookup versus new mathematics}

\begin{sloppypar}
Lookup: \code{Polynomial.coeff\_mul},
\code{Polynomial.finset\_sum\_coeff}, \code{Polynomial.coeff\_C\_mul},
\code{Polynomial.coeff\_X\_pow}, \code{Finset.sum\_ite\_eq},
\code{Finset.mem\_antidiagonal}, \code{pow\_le\_pow\_right}$_0$,
\code{one\_le\_pow}$_0$. New: only the growth-exponent bookkeeping.
\end{sloppypar}

\section{What the next tide inherits}

The architecture consult's implementation order for the rest of
stage 5: (1) the $T_2/H$ bridge and a named positive Gaussian rate;
(2) arbitrary-small Peano control on the mesoscopic window;
(3) absorption of the absolute corrections into a weakened Gaussian;
(5) the unrestricted exponential remainder
$|e^x - \sum_{i \le N} x^i/i!| \le |x|^{N+1} e^{|x|}/(N{+}1)!$ and
the perturbation bound $|e^{-(A+\delta)} - e^{-A}| \le |\delta|\,
e^{|A|+|\delta|}$; (6) the normalized $q$-uniform Gaussian majorant
(the single interface the dominated-convergence step consumes);
(7) filter-indexed DCT with the window indicator folded into the
integrand; (8) the two outer-tail removals via stage 2;
(9) \code{numerator\_hasExpansion}.
\end{document}
